%=====================================================================
%  Casos de Uso - Documento base (solo estructura)
%  Compilar con: pdflatex casos-de-uso.tex  (dos veces, por el indice)
%=====================================================================
\documentclass[11pt,letterpaper]{article}

%---------------------------------------------------------------------
% Paquetes
%---------------------------------------------------------------------
\usepackage[utf8]{inputenc}
\usepackage[T1]{fontenc}
\usepackage[spanish,es-tabla]{babel}
\usepackage[letterpaper,margin=2.5cm,includehead,headsep=8mm]{geometry}
\usepackage{array}
\usepackage{longtable}
\usepackage{enumitem}
\usepackage{fancyhdr}
\usepackage{lastpage}
\usepackage{ragged2e}
\usepackage[hidelinks]{hyperref}

%---------------------------------------------------------------------
% Datos del documento (editar aqui)
%---------------------------------------------------------------------
\newcommand{\tituloDoc}{Casos de Uso}
\newcommand{\sistema}{Bastion}
\newcommand{\versionDoc}{Versión 1.0}
\newcommand{\fechaDoc}{\today}
\newcommand{\organizacion}{Universidad Veracruzana}
\newcommand{\encabezadoCorto}{Casos de uso para Bastion}

%---------------------------------------------------------------------
% Encabezado y pie de pagina
%---------------------------------------------------------------------
\pagestyle{fancy}
\fancyhf{}
\fancyhead[L]{\footnotesize\itshape\bfseries \encabezadoCorto}
\fancyhead[R]{\footnotesize\itshape\bfseries Página \thepage}
\renewcommand{\headrulewidth}{0pt}

%---------------------------------------------------------------------
% Tipos de columna
%---------------------------------------------------------------------
\newcolumntype{L}[1]{>{\RaggedRight\arraybackslash}p{#1}}
\newcolumntype{J}[1]{>{\justifying\arraybackslash}p{#1}}

% Anchos de las columnas de la plantilla de caso de uso
\newlength{\colEtiqueta}\setlength{\colEtiqueta}{0.20\textwidth}
\newlength{\colValor}   \setlength{\colValor}{0.26\textwidth}
\newlength{\colEtiquetaB}\setlength{\colEtiquetaB}{0.21\textwidth}
\newlength{\colValorB}  \setlength{\colValorB}{0.21\textwidth}
% Ancho de una celda que ocupa las 3 ultimas columnas
\newlength{\colAncho}
\setlength{\colAncho}{\dimexpr\colValor+\colEtiquetaB+\colValorB+4\tabcolsep+2\arrayrulewidth\relax}

%---------------------------------------------------------------------
% Plantilla de caso de uso
%
%   \begin{casodeuso}{CU-00. Nombre del caso de uso}
%       \cuDoble{Creado Por:}{...}{Fecha de Creación:}{...}
%       \cuFila{Actor Primario:}{...}
%       ...
%   \end{casodeuso}
%---------------------------------------------------------------------
\newenvironment{casodeuso}[1]
{%
  \par\vspace{1em}\noindent
  \setlength{\extrarowheight}{2pt}%
  \begin{longtable}{|L{\colEtiqueta}|L{\colValor}|L{\colEtiquetaB}|L{\colValorB}|}
  \hline
  \textbf{ID y Nombre del Caso de Uso:} & \multicolumn{3}{l|}{\textbf{#1}} \\ \hline
  \endfirsthead
  \hline
  \multicolumn{4}{|l|}{\textit{#1 (continuación)}} \\ \hline
  \endhead
}
{%
  \end{longtable}\par\vspace{1em}%
}

% Fila simple: etiqueta + contenido a todo lo ancho
\newcommand{\cuFila}[2]{#1 & \multicolumn{3}{J{\colAncho}|}{#2} \\ \hline}

% Fila doble: dos pares etiqueta/valor en el mismo renglon
\newcommand{\cuDoble}[4]{#1 & #2 & #3 & #4 \\ \hline}

% Lista numerada compacta para flujos
\newenvironment{flujo}
{\begin{enumerate}[leftmargin=*,topsep=2pt,itemsep=1pt,parsep=0pt]}
{\end{enumerate}}

%=====================================================================
\begin{document}

%---------------------------------------------------------------------
% PORTADA
%---------------------------------------------------------------------
\thispagestyle{empty}
\noindent\rule{\textwidth}{2pt}
\begin{flushright}
  {\Huge\bfseries\sffamily \tituloDoc}\\[2.5em]
  {\large\bfseries\sffamily para}\\[2.5em]
  {\Huge\bfseries\sffamily \sistema}\\[2.5em]
  {\large\bfseries\sffamily \versionDoc}\\[2.5em]
  {\large\bfseries\sffamily Preparado por}\\[2.5em]
  {\large\bfseries\sffamily Autor}\\[2.5em]
  {\large\bfseries\sffamily \organizacion}\\[2.5em]
  {\large\bfseries\sffamily \fechaDoc}
\end{flushright}
\newpage

%---------------------------------------------------------------------
% SECCIÓN 1: HISTORIAL DE REVISIÓN
%---------------------------------------------------------------------
\section*{Historial de revisión}
\addcontentsline{toc}{section}{Historial de revisión}

\noindent
\begin{longtable}{|L{0.22\textwidth}|L{0.14\textwidth}|L{0.40\textwidth}|L{0.14\textwidth}|}
\hline
\textbf{Nombre} & \textbf{Fecha} & \textbf{Razón de cambios} & \textbf{Versión} \\ \hline
\endfirsthead
\hline
\textbf{Nombre} & \textbf{Fecha} & \textbf{Razón de cambios} & \textbf{Versión} \\ \hline
\endhead
  &   &   &   \\ \hline
  &   &   &   \\ \hline
\end{longtable}
\newpage

%---------------------------------------------------------------------
% SECCIÓN 2: LISTA DE CASOS DE USO
%---------------------------------------------------------------------
\section*{Lista de Casos de Uso}
\addcontentsline{toc}{section}{Lista de Casos de Uso}

\noindent
\begin{longtable}{|L{0.35\textwidth}|L{0.55\textwidth}|}
\hline
\textbf{\textit{Actor primario}} & \textbf{\textit{Caso de uso}} \\ \hline
\endfirsthead
\hline
\textbf{\textit{Actor primario}} & \textbf{\textit{Caso de uso}} \\ \hline
\endhead
  & CU-01  \\ \hline
  & CU-02  \\ \hline
  & CU-03  \\ \hline
\end{longtable}
\newpage

%---------------------------------------------------------------------
% SECCIÓN 3: CASOS DE USO
%---------------------------------------------------------------------
\section*{Casos de uso}
\addcontentsline{toc}{section}{Casos de uso}

%---------------------------- CU-01 ----------------------------------
\begin{casodeuso}{CU-01. Nombre del caso de uso}
  \cuDoble{Creado Por:}{}{Fecha de Creación:}{}
  \cuDoble{Actor Primario:}{}{Actores Secundarios:}{}
  \cuFila{Disparador:}{}
  \cuFila{Descripción:}{}
  \cuFila{Precondiciones:}{PRE-01: }
  \cuFila{Postcondiciones:}{POS-01: }
  \cuFila{Flujo Normal:}{%
    \begin{flujo}
      \item
      \item Termina el caso de uso
    \end{flujo}}
  \cuFila{Flujos Alternativos:}{FA-01:
    \begin{flujo}
      \item
    \end{flujo}}
  \cuFila{Excepciones:}{EX-01:
    \begin{flujo}
      \item
    \end{flujo}}
  \cuFila{Prioridad:}{}
  \cuFila{Frecuencia de Uso:}{}
  \cuFila{Reglas de Negocio:}{RN-}
  \cuFila{Otra Información:}{N/A}
  \cuFila{Suposiciones:}{N/A}
\end{casodeuso}

%---------------------------- CU-02 ----------------------------------
\begin{casodeuso}{CU-02. Nombre del caso de uso}
  \cuDoble{Creado Por:}{}{Fecha de Creación:}{}
  \cuDoble{Actor Primario:}{}{Actores Secundarios:}{}
  \cuFila{Disparador:}{}
  \cuFila{Descripción:}{}
  \cuFila{Precondiciones:}{PRE-01: }
  \cuFila{Postcondiciones:}{POS-01: }
  \cuFila{Flujo Normal:}{%
    \begin{flujo}
      \item
      \item Termina el caso de uso
    \end{flujo}}
  \cuFila{Flujos Alternativos:}{FA-01:
    \begin{flujo}
      \item
    \end{flujo}}
  \cuFila{Excepciones:}{EX-01:
    \begin{flujo}
      \item
    \end{flujo}}
  \cuFila{Prioridad:}{}
  \cuFila{Frecuencia de Uso:}{}
  \cuFila{Reglas de Negocio:}{RN-}
  \cuFila{Otra Información:}{N/A}
  \cuFila{Suposiciones:}{N/A}
\end{casodeuso}

% Copiar el bloque anterior para CU-03, CU-04, ... CU-NN

\newpage

%---------------------------------------------------------------------
% SECCIÓN 4: REGLAS DE NEGOCIO
%---------------------------------------------------------------------
\section*{Reglas de Negocio}
\addcontentsline{toc}{section}{Reglas de Negocio}

\noindent
\begin{longtable}{|L{0.20\textwidth}|J{0.45\textwidth}|L{0.22\textwidth}|}
\hline
\textbf{ID Regla} & \textbf{Regla del Negocio} & \textbf{Tipo} \\ \hline
\endfirsthead
\hline
\textbf{ID Regla} & \textbf{Regla del Negocio} & \textbf{Tipo} \\ \hline
\endhead
RN-01 &   & Restricción \\ \hline
RN-02 &   & Restricción \\ \hline
RN-03 &   & Restricción \\ \hline
\end{longtable}

\end{document}
